% =========================================================
% Breadcrumb
% =========================================================
\begin{tcolorbox}[
  colback=gray!6,
  colframe=gray!40,
  arc=2pt,
  left=8pt, right=8pt, top=6pt, bottom=6pt,
  title={\small\textbf{Where You Are in the Journey}},
  fonttitle=\small\bfseries
]
\begin{center}
\small
Sets \& Functions
$\;\to\;$ $\mathbb{R}$
$\;\to\;$ Algebraic Structures
$\;\to\;$ Abstract Algebra
$\;\to\;$ \textbf{Algebraic Geometry}
$\;\to\;$ $\cdots$
\end{center}

\medskip
\noindent\textbf{How we got here.}
Abstract algebra gave us the theory of polynomial rings and ideals.
Algebraic geometry asks the geometric question: what do the
\emph{solutions} of polynomial equations look like?
Every polynomial ideal determines a geometric object; every geometric
object corresponds to an ideal.

\medskip
\noindent\textbf{What this chapter builds.}
We study affine varieties (solution sets of polynomial systems),
the Nullstellensatz (the fundamental correspondence between ideals
and varieties), and Gröbner bases as a computational tool for
working with polynomial ideals.

\medskip
\noindent\textbf{Where this leads.}
Algebraic geometry in its modern form uses scheme theory and
cohomology. This chapter provides the classical foundations.
\end{tcolorbox}
\vspace{1em}

% =========================================================
% Algebraic Geometry — Structural Roadmap
% Driver Text: Emily Clader & Dustin Ross, Beginning in Algebraic Geometry
% File: algebraic-geometry/algebraic-geometry.tex
% =========================================================
\section{Algebraic Geometry}
% =========================================================
% Structural Roadmap
% =========================================================
\section*{Structural Roadmap}
The development of algebraic geometry in this project follows
the definition--theorem--structure architecture used
throughout the analysis volumes.
The primary driver is \textit{Beginning in Algebraic Geometry}
by Emily Clader and Dustin Ross. The emphasis is on the interplay
between polynomial algebra and geometric structure,
building from affine varieties through to projective geometry
and culminating topics in modern algebraic geometry.
Each major topic is organized as:
\begin{center}
\textbf{Definitions $\longrightarrow$ Main Theorems $\longrightarrow$ Consequences and Structural Insight}
\end{center}
The global progression follows Clader--Ross in three stages:
\begin{enumerate}
  \item \textbf{Algebraic Foundations}
  \begin{enumerate}
    \item Polynomial rings
  \end{enumerate}
  \item \textbf{Affine Algebraic Geometry}
  \begin{enumerate}
    \item Varieties and ideals
    \item Irreducibility of affine varieties
    \item Coordinate rings
    \item Polynomial maps
    \item Proof of the Nullstellensatz
    \item Dimension
    \item Smoothness
    \item Products
  \end{enumerate}
  \item \textbf{Projective Algebraic Geometry}
  \begin{enumerate}
    \item Projective varieties
    \item Maps of projective varieties
    \item Quasiprojective varieties
    \item Culminating topics
  \end{enumerate}
\end{enumerate}
\vspace{1em}
\begin{remark}
This treatment prioritizes geometric intuition grounded in
rigorous commutative algebra. The book opens with polynomial rings
as the algebraic foundation before any geometric objects are introduced,
reflecting the philosophy that the algebra must be solid before
the geometry can be understood.
\end{remark}
\begin{remark}
The central objects of study are varieties and the maps between them.
The Nullstellensatz is the pivotal theorem of the affine theory,
establishing the precise correspondence between geometry and algebra
that underlies all subsequent development.
\end{remark}
\begin{remark}[Structural Position]
Algebraic geometry is developed building on the commutative algebra
and linear algebra established in earlier chapters.
The affine theory is developed completely before projective geometry
is introduced, so that projective space can be understood as a
natural compactification and generalization of affine space.
\end{remark}
% =========================================================
% 0. Algebraic Foundations
% =========================================================
 \input{volume-iv/algebra/algebraic-geometry/notes/notes-polynomial-rings}
% =========================================================
% I. Affine Algebraic Geometry
% =========================================================
% \input{volume-iv/algebra/algebraic-geometry/notes/notes-varieties-ideals}
% \input{volume-iv/algebra/algebraic-geometry/notes/notes-irreducibility}
% \input{volume-iv/algebra/algebraic-geometry/notes/notes-coordinate-rings}
% \input{volume-iv/algebra/algebraic-geometry/notes/notes-polynomial-maps}
% \input{volume-iv/algebra/algebraic-geometry/notes/notes-nullstellensatz}
% \input{volume-iv/algebra/algebraic-geometry/notes/notes-dimension}
% \input{volume-iv/algebra/algebraic-geometry/notes/notes-smoothness}
% \input{volume-iv/algebra/algebraic-geometry/notes/notes-products}
% =========================================================
% II. Projective Algebraic Geometry
% =========================================================
% \input{volume-iv/algebra/algebraic-geometry/notes/notes-projective-varieties}
% \input{volume-iv/algebra/algebraic-geometry/notes/notes-maps-projective}
% \input{volume-iv/algebra/algebraic-geometry/notes/notes-quasiprojective}
% \input{volume-iv/algebra/algebraic-geometry/notes/notes-culminating-topics}
% =========================================================
% Future Sections (To Be Developed)
% =========================================================
% \input{volume-iv/algebra/algebraic-geometry/notes/notes-schemes}
% \input{volume-iv/algebra/algebraic-geometry/notes/notes-sheaves}
% \input{volume-iv/algebra/algebraic-geometry/notes/notes-cohomology}
% =========================================================
% Logical Positioning
% =========================================================